\section{Half-integer orders and the Euclidean chain}\label{inf:sec-parity}

The order parameter belongs to $\tfrac12\Z$ at a Euclidean input.
The block dimensions $a,b$ and the angular and radial degree indices
are integers; the Jacobi parameters and the Bessel orders may be
half-integers.

\begin{lemma}[Half-integer non-coincidence]\label{inf:half-integer-noncoincidence}
For $\nu\in\{1/2,3/2,\ldots\}$, $J_\nu$ has no positive algebraic
zero. Two distinct positive half-integer orders have no common positive
zero. In particular $J_p$ and $J_{p+3}$ have no common positive zero
for $p\in\Z+1/2$, $p>0$.
\end{lemma}
\begin{proof}
Write $\nu=l+1/2$, $l\ge0$. The half-order identities and recurrence
\cite[\S\S10.6, 10.16]{inf:dlmf} give
\[
 \sqrt{\pi x/2}\,J_{l+1/2}(x)
 =\mathsf F_l(x)\sin x+\mathsf G_l(x)\cos x,
\]
where $\mathsf F_l,\mathsf G_l\in\mathbb Q[x^{-1}]$ and
\[
 (\mathsf F_{-1},\mathsf G_{-1})=(0,1),\qquad
 (\mathsf F_0,\mathsf G_0)=(1,0),\qquad
 (\mathsf F_{l+1},\mathsf G_{l+1})
 =\frac{2l+1}{x}(\mathsf F_l,\mathsf G_l)
       -(\mathsf F_{l-1},\mathsf G_{l-1}).
\]
The determinant of the successive rows $(l,l-1)$ is one: the
recurrence preserves it and the initial determinant is one.
Thus the two entries in row $l$ never vanish together at $x\ne0$.
If $x>0$ were algebraic and a zero, the real algebraic numbers
$\mathsf F_l(x),\mathsf G_l(x)$ would satisfy
\[
 e^{2ix}=-\frac{\mathsf G_l(x)+i\mathsf F_l(x)}
                  {\mathsf G_l(x)-i\mathsf F_l(x)}.
\]
The denominator is nonzero. This contradicts the Hermite--Lindemann
theorem, which states that $e^z$ is transcendental for each nonzero
algebraic complex number $z$ \cite[\S1.2]{inf:BakerWustholz}.

At a zero $x$ of $J_\nu$, $J_{\nu+1}(x)\ne0$: otherwise the derivative
recurrence and uniqueness for Bessel's equation give $J_\nu=0$
identically. Define the Lommel polynomials by
\[
 \mathsf L^{\mathrm{Lom}}_{0,\mu}(x)=1,\qquad
 \mathsf L^{\mathrm{Lom}}_{1,\mu}(x)=2\mu/x,\qquad
 \mathsf L^{\mathrm{Lom}}_{k+1,\mu}(x)
 =\frac{2(\mu+k)}x\mathsf L^{\mathrm{Lom}}_{k,\mu}(x)
     -\mathsf L^{\mathrm{Lom}}_{k-1,\mu}(x).
\]
The Bessel recurrence implies
$J_{\nu+k}(x)=\mathsf L^{\mathrm{Lom}}_{k-1,\nu+1}(x)J_{\nu+1}(x)$
for $k\ge1$; see \cite{inf:watson}.
For $k\ge2$ this is a nonzero polynomial in $x^{-1}$ with rational
coefficients and leading coefficient $\prod_{h=1}^{k-1}2(\nu+h)$.
A common zero would therefore make $x$ algebraic, contrary to the
first assertion. For $k=1$ the non-coincidence was already proved.
For the primary pair the identity specializes to
\[
 J_{p+3}(x)=\left\{\frac{4(p+1)(p+2)}{x^2}-1\right\}J_{p+1}(x)
 \quad\text{when }J_p(x)=0.
\]
\end{proof}

\begin{lemma}[Shifted congruence approximation]\label{inf:shifted-arithmetic}
The set $\mathcal P_{\mathrm{bd}}$ has infinitely many members in each
of $\Z$ and $\Z+1/2$. For the pairs produced by
Lemma~\ref{inf:congruence} applied to $\alpha_{\mathrm{ar}}+1$, set
$\widetilde P=P-Q$ and $p=(\widetilde P-3)/2$. Along this particular
half-integer selection,
\[
 \limsup\widehat p_{(Q-1)/4}|p-\widehat p_{(Q-1)/4}|
 \le\frac98\alpha_{\mathrm{ar}}+d_{\mathrm{ar}}<31/10.
\]
\end{lemma}
\begin{proof}
Both $P,Q$ are odd, so $\widetilde P$ is even and $2p$ is odd.
Moreover $Q\equiv1\pmod4$ and
\[
 Q|Q\alpha_{\mathrm{ar}}-\widetilde P|
 =Q|Q(\alpha_{\mathrm{ar}}+1)-P|<9/2.
\]
The denominators tend to infinity. Put $m=(Q-1)/4$ and use
\eqref{inf:root-expansion}; multiplication by $\widehat p_m$ gives
\[
 \widehat p_m|p-\widehat p_m|
 \le\tfrac12\widehat p_m|Q\alpha_{\mathrm{ar}}-\widetilde P|
       +d_{\mathrm{ar}}\widehat p_m/A_m+O(A_m^{-1}).
\]
Now $\widehat p_m/Q\to\alpha_{\mathrm{ar}}/2$ and
$\widehat p_m/A_m\to1$. The rational comparison in
\eqref{inf:arithmetic-limsup} proves the bound. Positivity of $p$ and
$m\ge m_0$ hold eventually. Lemma~\ref{inf:half-integer-noncoincidence}
excludes zero distance. The integer selection is
Theorem~\ref{inf:near-integer}.
\end{proof}

\begin{proposition}[The half-lattice analytic chain]\label{inf:parity-chain}
The analytic statements listed in the table below that use a
near-primary input hold for
$p\in\mathcal P_{\mathrm{bd}}$, uniformly for $0<\tau\le20$ and
$r\in[1/4,3/8]$, with integer $a,b$ for each Euclidean conclusion.
The remaining analytic statements hold under their displayed real
parameter hypotheses.
\end{proposition}
\begin{proof}
The chain is given below, with its integer objects and the identities
that determine the permissible parameters. Each row is applied under
the hypotheses of its indicated statements.
\begingroup\small\normalfont
\renewcommand{\arraystretch}{1.12}
\begin{longtable}{@{}>{\raggedright\arraybackslash}p{.32\textwidth}>{\raggedright\arraybackslash}p{.62\textwidth}@{}}
\toprule Statements & Parameter dependence\\\midrule
\endfirsthead
\toprule Statements & Parameter dependence\\\midrule
\endhead
Lemma~\ref{inf:reduction} &
The coordinate Laplacian and beta measure use integer $a,b$ and
$p=(a+b)/2$. The invariant angular multiplicity is one.\\
Lemma~\ref{inf:angular-formulas} &
Rodrigues integration and rising factorials use positive real
$\ab,\bb$; $j$ is the integer polynomial degree.\\
Lemma~\ref{inf:angular-sup} &
The endpoint comparison needs $\bb-1\ge\ab-1\ge0$.
Bernstein applies on great circles in every integer dimension.\\
Lemma~\ref{inf:solid-harmonics} &
$t^jG_j$ is a polynomial in $q_X,q_Y$.
The regular Bessel mode is this polynomial times an entire power
series in $\varrho^2$, including the origin and axes.\\
Proposition~\ref{inf:ball-spectrum} &
The radial order $p-1+2j>1$ is real; its singular solution is
outside the Friedrichs form domain. Integer blocks identify the
Euclidean invariant realization.\\
Lemma~\ref{inf:radial-polynomials} &
The radial parameter $2j+p-1$ is real. Rodrigues and generalized
binomial identities require only integer polynomial indices.\\
Lemmas~\ref{inf:bessel-spacing}, \ref{inf:bessel-order} &
The radial ODE, Poincar\'e inequality and Watson integral have
their stated real-order hypotheses.\\
Lemmas~\ref{inf:debye-first}, \ref{inf:debye-finite} &
The outgoing Volterra equation is for real $\nu>0$;
only its finite expansion order is integral.\\
Lemma~\ref{inf:window-isolation} &
The recurrence shifts the real Bessel order by two.
Its frequency spacing concerns all radial indices.\\
Lemma~\ref{inf:phase-strip}, Theorem~\ref{inf:global-count} &
The lattice coordinates are the integer angular index and
integer phase index; their offsets may depend on real $p$.\\
Lemma~\ref{inf:congruence}, Theorem~\ref{inf:near-integer},
Lemma~\ref{inf:shifted-arithmetic} &
Apply the same irrational approximation to $\alpha_{\mathrm{ar}}$
or $\alpha_{\mathrm{ar}}+1$. The latter changes $P$ to $P-Q$
and preserves the error product.\\
Proposition~\ref{inf:primary-expansion} &
The root expansion and uniqueness use real-order continuation.
Bourget and Lemma~\ref{inf:half-integer-noncoincidence}
exclude both half-lattice coincidences.\\
Lemma~\ref{inf:detuning} &
The exact rational trace identity and the real-order derivative
give the same multiplicative detuning relation.\\
Lemmas~\ref{inf:riccati-gap}, \ref{inf:low-recurrence} &
The finite integer order shifts leave the base order real.
The period-six sequence comes from the limiting recurrence.\\
Proposition~\ref{inf:low-exclusion} &
All relevant zeros are continued on a real interval in their
orders; the selected endpoint may be a half-integer.\\
Lemma~\ref{inf:angular-weighted} &
The entry and derivative formulas are rational identities for
positive real beta parameters with sum $p$.\\
Lemma~\ref{inf:angular-positive} &
The coefficient certificates hold for all real $p\ge64$.
The recurrence-ratio bounds use $\ab,\bb\ge16$.\\
Proposition~\ref{inf:flow-form} &
Axis cutoffs use $\ab,\bb>1$ and the origin cutoff $p>1$.
Integer blocks identify the weighted and Euclidean closures.\\
Lemmas~\ref{inf:flow-relative}, \ref{inf:transport} &
The band estimates and Green identity use real Jacobi parameters
and radial Bessel norms; both angular parities are retained.\\
Lemma~\ref{inf:flow-complement} &
The complementary inverse is radial spectral calculus and a
relative-form Neumann series on every complementary radial state.\\
Lemma~\ref{inf:flow-low-block} &
The finite recurrence gives seed offset $-\theta/(2p)$ and
index-five offset $-72$ in either ambient parity.\\
Lemma~\ref{inf:ordered-eigenvalues}, Theorem~\ref{inf:flow-theorem} &
Finite Riesz reduction and the weighted derivative margin apply
on the common real-parameter form domain, in the stated narrow band.\\
Theorem~\ref{inf:integer-gap} &
Integers in $[p/2,3p/4]$ number at least $p/4-1$ and have
$\varsigma$ spacing at least $1/(8p)$.
$2p\in\Z$ ensures integer $b=2p-a$.\\
Theorem~\ref{inf:centres} &
The finite recursion uses $\eps=1/p$, fixed degree indices,
the divisor $-1/8$, positive $c_j$ and powers of $r(1-r)$.\\
Lemma~\ref{inf:analytic-coefficients},
Corollary~\ref{inf:centre-coefficient-defect} &
The analytic coefficient transfer uses any positive measure and
$0<\mathsf H_j(1)=(\bb)_j/(p+j-1)_j\le1$.\\
Proposition~\ref{inf:centre-gap-transfer} &
Indexwise min--max and radial dilation hold in every integer dimension;
the normalized error is $O_N(\tau p^{-4})$.\\
Definition~\ref{inf:spaces}, Lemma~\ref{inf:radial-stencils},
Proposition~\ref{inf:module} &
The grade exponent is real. Gasper uses real Jacobi parameters;
the integer raising/lowering shifts are angular degree differences.\\
Lemma~\ref{inf:poisson-columns}, Proposition~\ref{inf:poisson-bounds} &
The radial columns are rational in the real parameter $\beta$;
the separate harmonic source column remains included.\\
Lemmas~\ref{inf:raising-coupling}, \ref{inf:clamped-smoothing} &
Positive radial stencils and their coupling use real $\beta\ge0$;
the smoothing gain acts on the clamped layers $s\ge1$.\\
Lemma~\ref{inf:gradient-profiles},
Proposition~\ref{inf:grouped-shapes} &
The solid product identity uses integer solid degrees and real $p$;
the grouped high-shape estimates have clamped source domains.\\
Lemmas~\ref{inf:double-lift}, \ref{inf:derivative-moments} &
The two lift columns and derivative columns use real radial
parameters; their reference normal factor is $R=O(p)$.\\
Lemma~\ref{inf:radial-centre-norm},
Proposition~\ref{inf:centre-source-norm} &
Rodrigues--Sonine, Lommel--Parseval and the positive Bessel ratio
hold for real orders. The head cutoff satisfies $\nu+1+2S\ge2R$;
composition uses the finite list of known centre modes.\\
Lemma~\ref{inf:pullback}, Proposition~\ref{inf:nonlinear-map} &
The Jacobian is $q^{2p-1}d=q^{n-1}d$ at integer $n$.
The nonlinear map has clamped source domain $\Gc_p\times\Ec_1$.\\
Proposition~\ref{inf:residual-estimates},
Theorem~\ref{inf:material-isomorphism} &
The double lift, normal Hessian and material identities use the
same real parameters; the $L^2$ completion is at fixed integer $n$.\\
Lemmas~\ref{inf:mixed-algebra}, \ref{inf:mixed-Poisson},
\ref{inf:mixed-domain}, Theorem~\ref{inf:relative-graph} &
The radial Hilbert blocks use real weights; density uses real
binomial powers. Hilbert and coefficient estimates retain their
respective graph domains and all radial layers.\\
Theorem~\ref{inf:local-count}, Corollary~\ref{inf:clusters} &
The integer phase lattice and real turning-point ODE give the
same local count and cluster cardinality.\\
Lemma~\ref{inf:cluster-inverse},
Proposition~\ref{inf:mixed-inverse} &
The cofactor argument is finite linear algebra, followed by the
radial graph reconstruction in both spaces.\\
Lemma~\ref{inf:head-tail}, Theorem~\ref{inf:coefficient-inverse} &
The head has at most $K_{\mathrm{head}}p+1$ radial indices;
the infinite tail uses the same all-layer Poisson columns.\\
Theorem~\ref{inf:newton-threshold}, Corollary~\ref{inf:Newton-zero} &
The contraction uses the displayed powers of $p,\tau$ in the same
real Banach spaces; its centre order is fixed first.\\
Lemmas~\ref{inf:geometry-convexity}, \ref{inf:geometry-nonball},
Propositions~\ref{inf:classical-solution}, \ref{inf:Pompeiu} &
The sphere and elliptic arguments require integer $n,a,b$.
The regular solid expansion covers the axes and origin in each parity.\\
\bottomrule
\end{longtable}
\endgroup
The only lattice changes are thus the non-coincidence theorem,
the arithmetic choice of $p$, the integer split grid, and the
integer block realization. All estimates in the table retain their
stated constants on this half-lattice.
\end{proof}
